\documentclass{report}
\usepackage{mathesis}
\usepackage[utf8]{inputenc}
\usepackage[russian]{babel}

\begin{document}

\section{Класс интегрируемых по Риману функций (Riemann Integrable Class)}\label{entity:obj-riemann-class}
\textbf{Тип:} object \quad \textbf{Источник:} Зорич В.А., Математический анализ, Том I (zorich-1, p. 1)

\textbf{Каноническая запись:}
\[
 \mathcal{R}\entityref{obj-closed-interval}{[a,b]} \mDefIff \{ \entityref{obj-function}{f} \mid \mExists \int_a^b \entityref{obj-function}{f}(x) dx \} 
\]

\vspace{1em}
\hrule
\vspace{1em}

\section{Функция (Function)}\label{entity:obj-function}
\textbf{Тип:} object \quad \textbf{Источник:} Зорич В.А., Математический анализ, Том I (zorich-1, p. 1)

\textbf{Каноническая запись:}
\[
 f \subset \entityref{obj-set}{X} \times \entityref{obj-set}{Y} 
\]
\[
 \mForall{x} (x \in X \mImplies \mExists{! y} (y \in Y \land (x, y) \in f)) 
\]

\textbf{Естественный язык:}
Функция $f$ из множества $X$ в множество $Y$ — это правило (или бинарное отношение), по которому каждому элементу $x \in X$ ставится в соответствие ровно один элемент $y \in Y$. Формально, это подмножество декартова произведения $X \times Y$, обладающее свойством однозначности.

\vspace{1em}
\hrule
\vspace{1em}

\section{Вещественные числа (Real Numbers)}\label{entity:obj-real-numbers}
\textbf{Тип:} object \quad \textbf{Источник:} Зорич В.А., Математический анализ, Том I (zorich-1, p. 1)

\textbf{Каноническая запись:}
\[
 \mathbb{R} = (\entityref{obj-set}{R}, +, \cdot, \leq) 
\]
\[
 \text{Где } \mathbb{R} \text{ удовлетворяет аксиомам полного упорядоченного поля.} 
\]

\textbf{Естественный язык:}
Множество вещественных чисел $\mathbb{R}$ — это непрерывная числовая прямая. Формально $\mathbb{R}$ задается как полное (непрерывное) архимедово упорядоченное поле. В нем можно складывать, умножать, сравнивать элементы, и в нем нет «дырок» (каждое ограниченное сверху подмножество имеет точную верхнюю грань).

\vspace{1em}
\hrule
\vspace{1em}

\section{Замкнутый отрезок (Closed Interval)}\label{entity:obj-closed-interval}
\textbf{Тип:} object \quad \textbf{Источник:} Зорич В.А., Математический анализ, Том I (zorich-1, p. 1)

\textbf{Каноническая запись:}
\[
 [a,b] \mDefIff \{ x \in \entityref{obj-real-numbers}{\mathbb{R}} \mid a \le x \le b \} 
\]

\vspace{1em}
\hrule
\vspace{1em}

\section{Множество (Set)}\label{entity:obj-set}
\textbf{Тип:} object \quad \textbf{Источник:} Зорич В.А., Математический анализ, Том I (zorich-1, p. 1)

\textbf{Каноническая запись:}
\[
 \entityref{obj-set}{X} \text{ - базовое понятие ZFC, определяемое аксиомами:} 
\]
\[
 \entityref{axm-zfc-extensionality}{Axiom of Extensionality}, \entityref{axm-zfc-pairing}{Axiom of Pairing}, \entityref{axm-zfc-union}{Axiom of Union}, \entityref{axm-zfc-power-set}{Axiom of Power Set}, \entityref{axm-zfc-infinity}{Axiom of Infinity}, \entityref{axm-zfc-regularity}{Axiom of Regularity}, \entityref{axm-zfc-specification}{Axiom of Specification}, \entityref{axm-zfc-replacement}{Axiom of Replacement}, \entityref{axm-zfc-choice}{Axiom of Choice} 
\]

\textbf{Естественный язык:}
Множество — базовое, неопределяемое напрямую понятие математики. Множество представляет собой совокупность объектов произвольной природы, называемых его элементами. Все свойства множеств строго выводятся из аксиом системы Цермело-Френкеля (ZFC).

\vspace{1em}
\hrule
\vspace{1em}

\section{Axiom of Union}\label{entity:axm-zfc-union}
\textbf{Тип:} axiom \quad \textbf{Источник:} Unknown (Unknown)

\textbf{Каноническая запись:}
\[

\mForall{\mathcal{F}}
\mExists{U} \mForall{x} \left( x \in U \mIff \mExists{A} (A \in \mathcal{F} \land x \in A) \right)

\]

\vspace{1em}
\hrule
\vspace{1em}

\section{Axiom of Pairing}\label{entity:axm-zfc-pairing}
\textbf{Тип:} axiom \quad \textbf{Источник:} Unknown (Unknown)

\textbf{Каноническая запись:}
\[

\mForall{A, B}
\mExists{C} \mForall{x} \left( x \in C \mIff (x = A \lor x = B) \right)

\]

\vspace{1em}
\hrule
\vspace{1em}

\section{Axiom of Power Set}\label{entity:axm-zfc-power-set}
\textbf{Тип:} axiom \quad \textbf{Источник:} Unknown (Unknown)

\textbf{Каноническая запись:}
\[

\mForall{X}
\mExists{P} \mForall{A} (A \in P \mIff A \subset X)

\]

\vspace{1em}
\hrule
\vspace{1em}

\section{Axiom of Regularity}\label{entity:axm-zfc-regularity}
\textbf{Тип:} axiom \quad \textbf{Источник:} Unknown (Unknown)

\textbf{Каноническая запись:}
\[

\mForall{S}
\left( (\mExists{y} y \in S) \mImplies \mExists{x} \left( x \in S \land \lnot \mExists{z} (z \in S \land z \in x) \right) \right)

\]

\vspace{1em}
\hrule
\vspace{1em}

\section{Axiom of Infinity}\label{entity:axm-zfc-infinity}
\textbf{Тип:} axiom \quad \textbf{Источник:} Unknown (Unknown)

\textbf{Каноническая запись:}
\[

\mExists{I}
\left( \emptyset \in I \;\land\; \mForall{x} \left( x \in I \mImplies \mExists{y} (y \in I \land \mForall{z} (z \in y \mIff z \in x \lor z = x)) \right) \right)

\]

\vspace{1em}
\hrule
\vspace{1em}

\section{Axiom of Specification}\label{entity:axm-zfc-specification}
\textbf{Тип:} axiom \quad \textbf{Источник:} Unknown (Unknown)

\textbf{Каноническая запись:}
\[

\mForall{X}
\mExists{Y} \mForall{z} \left( z \in Y \mIff (z \in X \land \phi(z)) \right)

\]

\vspace{1em}
\hrule
\vspace{1em}

\section{Axiom of Choice}\label{entity:axm-zfc-choice}
\textbf{Тип:} axiom \quad \textbf{Источник:} Unknown (Unknown)

\textbf{Каноническая запись:}
\[

\mForall{\mathcal{F}}
\left(
  \mForall{A, B} \left( A \in \mathcal{F} \land B \in \mathcal{F} \land A \neq B \mImplies \lnot \mExists{x} (x \in A \land x \in B) \right)
  \land \lnot(\emptyset \in \mathcal{F})
\right)
\mImplies
\mExists{C} \mForall{A} \left( A \in \mathcal{F} \mImplies \mExists{! x} (x \in A \land x \in C) \right)

\]

\vspace{1em}
\hrule
\vspace{1em}

\section{Axiom of Extensionality}\label{entity:axm-zfc-extensionality}
\textbf{Тип:} axiom \quad \textbf{Источник:} Unknown (Unknown)

\textbf{Каноническая запись:}
\[

\mForall{A, B}
(A = B) \mIff (\mForall{x} x \in A \mIff x \in B)

\]

\vspace{1em}
\hrule
\vspace{1em}

\section{Axiom of Replacement}\label{entity:axm-zfc-replacement}
\textbf{Тип:} axiom \quad \textbf{Источник:} Unknown (Unknown)

\textbf{Каноническая запись:}
\[

\left( \mForall{x} \mExists{! y} \, \phi(x, y) \right)
\mImplies
\mForall{X} \mExists{Y} \mForall{y} \left( y \in Y \mIff \mExists{x} \left( x \in X \land \phi(x, y) \right) \right)

\]

\vspace{1em}
\hrule
\vspace{1em}



\end{document}
